\section{Introduction}\label{sec:introduction}

The primary failure mode in recursive multi-agent systems is not malfunction. It is \emph{drift}: the gradual, often undetected separation of a spawned agent's behavior from the intent of its originating principal, compounding silently across generations of recursive spawning until the system's observable trajectory bears no recognizable relation to any human-directed purpose. Drift does not require a defective or malicious agent. It emerges from the structure of delegation chains themselves — specifically, from the absence of a structural constraint that prevents intermediate nodes from severing or obscuring the causal link between a leaf agent's actions and the human principal who authorized its ancestry.

In uncontrolled recursive deployments, each spawn event creates a new accountability boundary. The parent agent — not the human principal — becomes the immediate authority over the child. If the parent is misaligned, operating under partial observability, or simply optimizing for a proxy objective, its delegation carries that uncertainty forward. The child inherits not only a task but an accountability posture: implicit assumptions about objectives, constraints, and escalation paths that may never have been reviewed by a human. When that child spawns its own descendants, the compounding is structural. By the third or fourth generation, the human principal may be factually excluded from the causal chain — not by deliberate exclusion, but by the accumulated weight of delegation chains that no layer was required to anchor to a named person.

\textbf{The Autonomous Drift Problem.} We call this the autonomous drift problem: in any multi-agent system that permits recursive spawning without structural accountability anchoring, the probability that a leaf agent's actions remain traceable to and controllable by a specific human principal decreases as a function of chain depth, and this decrease is not detected by the system because no layer maintains a verifiable chain of accountability. The problem is architectural, not behavioral. It cannot be solved by better prompting, improved alignment training, or stricter capability filtering because all of these solutions operate within the accountability boundary of the parent — and the parent's accountability boundary is precisely what drift erodes.

The field has attempted to address drift through two families of shallow-restriction mechanisms. The first is \textbf{time-to-live} (TTL) enforcement: spawn chains are permitted to extend only to a fixed depth or wall-clock limit, after which further spawning is blocked. The second is \textbf{depth limits} with static capability caps: agents above a given generation threshold lose access to certain tool classes, cannot write to certain state partitions, or require additional human confirmation before acting. Both are architecturally reasonable as circuit breakers and are insufficient as accountability mechanisms. TTL enforcement does not distinguish between a spawn chain that has correctly executed its purpose and one that has diverged from it — both terminate at the same depth. Depth limits treat capability restriction as a proxy for accountability, but a chain of many low-capability actions can collectively produce a high-capability outcome, and neither approach addresses the root structural condition: the absence of an immutable, verifiable chain connecting every node's authorization state to a named human principal.

\textbf{Core Insight.} The core structural insight of this paper is that drift is made structurally impossible by one mechanism alone: an \textbf{immutable parent pointer chain} enforced by the \Fact{} layer at every node, such that no agent — including the node's own parent — can modify, truncate, or obscure the chain of accountable ancestry linking any node to its originating human principal. The immutability is not a policy constraint applied to agent behavior — it is a structural property of the \Fact{} layer's write protocol. When a child node is spawned, the parent's \Fact{} layer atomically writes the child's identity and the parent's own \Fact{} layer reference into the child's \Fact{} state. From that moment forward, the child can always trace its ancestry by following parent pointers to their terminus. The terminus is always a node whose parent pointer is $\bot$ — the root node — and the root node's \Fact{} layer contains the identity of the human principal who provisioned it. No agent in the chain can alter this chain without being detected, because the chain is the evidence of its own integrity.

\textbf{Formal Definition of Drift.} We define drift formally as follows:

\begin{invariant}[Drift Condition]\label{invariant:drift}
A \bxthree{} node $n$ is in a state of \emph{drift} at time $t$ if and only if there exists a causal chain of agent actions $\alpha_1, \alpha_2, \dots, \alpha_k$ such that:
\begin{enumerate}
  \item $\alpha_1$ was authorized by a human principal $h$;
  \item each subsequent $\alpha_{i+1}$ was taken by a node spawned (directly or transitively) as a descendant of the node that took $\alpha_i$; and
  \item the \Fact{} layer of the node taking $\alpha_k$ does not contain a verifiable parent pointer chain traceable to $h$.
\end{enumerate}
\end{invariant}

Under this definition, drift is a structural condition, not a behavioral one. A node may exhibit drift without any individual agent acting incorrectly. Detection operates by inspecting the integrity of the parent pointer chain maintained in the \Fact{} layer — not by monitoring behavioral outputs. The \bxthree{} Recursive Spawning pillar guarantees:

\begin{property}[Drift Impossibility]\label{prop:drift-impossibility}
In a \bxthree{} system in which the Recursive Spawning pillar is correctly implemented, no node can exist in the state defined in Invariant~\ref{invariant:drift}.
\end{property}

This guarantee holds because the immutable parent pointer chain is written atomically at spawn time by the \Fact{} layer and is non-repudiable: the parent cannot modify, revoke, or obscure the pointer after the child has been instantiated. The child's \Fact{} layer validates its own ancestry on every decision cycle. If the chain is broken or untraceable to a named human principal, the node triggers the Bailout Protocol immediately, without discretionary routing through the parent.

\textbf{Roadmap.} The remainder of this paper is organized as follows. Section~\ref{sec:background} situates the Recursive Spawning pillar within the broader literature on multi-agent accountability and recursive system safety. Section~\ref{sec:three-layer} formally introduces the \bxthree{} three-layer model as the architectural substrate. Section~\ref{sec:recursive-spawning-pillar} provides the full formal specification of the Recursive Spawning pillar, including the immutable parent pointer chain mechanism, the spawn authorization protocol, and the ancestry verification routine executed by the \Fact{} layer on each decision cycle. Section~\ref{sec:pillar-interactions} addresses the interactions between Recursive Spawning and the other four \bxthree{} pillars: Loop Isolation, Spatial Firewall, Sandbox Gate, and Bailout Protocol. Section~\ref{sec:correctness} establishes the correctness guarantees, including the drift impossibility theorem and characterization of pillar behavior under adversarial spawn conditions. Section~\ref{sec:related} evaluates the Recursive Spawning approach against existing spawning governance models including TTL enforcement, capability-tiered depth limits, and dynamic delegation frameworks. Section~\ref{sec:limitations} addresses structural limitations and the directed research agenda. Section~\ref{sec:conclusion} concludes.
